\documentclass[11pt]{article}
\usepackage[margin=1in]{geometry}
\usepackage{lmodern}
\usepackage{hyperref}
\usepackage{enumitem}
\usepackage{amsmath,amssymb,amsthm,mathtools}
\usepackage[T1]{fontenc}
\usepackage[utf8]{inputenc}
\title{\textbf{Theory-mode report: A transition in the hole probability at finite temperature for free fermions in \$d\$ dimensions}}
\author{Generated from Scholar Inbox digest}
\date{Digest date: 2026-05-08}
\begin{document}
\maketitle
\section*{Paper information}
\begin{itemize}[leftmargin=1.5em]
\item \textbf{Title:} A transition in the hole probability at finite temperature for free fermions in \$d\$ dimensions
\item \textbf{Authors:} Giuseppe Del Vecchio Del Vecchio, Pierre Le Doussal, Gregory Schehr
\item \textbf{arXiv:} \href{https://arxiv.org/abs/2605.05431}{2605.05431}
\item \textbf{Scholar digest score:} 94.0
\end{itemize}
\section*{Executive summary}
This report summarizes the highest-scoring paper in the Scholar Inbox digest dated \textbf{2026-05-08}. The abstract states the core claim as follows: \emph{In a free Fermi gas at temperature \$T\$ much higher than the Fermi temperature one expects that the fluctuations of the number of particles in a given region has Poissonian/classical statistics. On the other hand at low temperature the Pauli exclusion principle leads to non trivial counting statistics. It is of great interest from a theoretical and experimental point of view to characterize the crossover between these two limits. Here we focus on the hole probability \$P(R,T)\$, i.e. the probability that a region of size \$R\$ is devoid of particles, in dimension \$d\$, and on the case of a spherical region of large radius \$R\$. We show that at low temperature it takes the scaling form \$P(R,T)\textbackslash\{\}sim \textbackslash\{\}exp\textbackslash\{\}big[-(k\_F R)\textasciicircum{}\{d+1\}Φ\_d(u=2R\textbackslash\{\},T/k\_F)\textbackslash\{\}big],\$ where \$k\_F\$ is the Fermi momentum. By mapping the problem to an effective Coulomb gas, we compute exactly the scaling function \$Φ\_d(u)\$ in any dimension. }.

\section*{Problem setup and intuition}
The introduction reinforces this lens: the paper appears motivated by a gap between practical behavior and theoretical understanding.

\section*{Proof architecture}
The technical core appears to build the surrogate quantity and the chain of lemmas needed to propagate local control into the final theorem: the paper follows a standard theory-paper structure with structural reduction plus quantitative bounds.

\section*{Takeaway}
The key reusable lesson is to define the right object, find the structural reduction, and quantify the error terms clearly.
\end{document}
